% In this chapter, we have reviewed basic skills for finding indefinite integrals. We've looked at the antiderivative formulas for all of the basic functions and reviewed techniques for finding antiderivatives of other functions.
%
% We've also reviewed the more advanced techniques of integration by partial fractions and integration by parts, both topics only for the BC Calculus course .

\question
\[
\displaystyle\int\left(3 x^{2}-2 x+3\right) \,dx=
\]
\begin{choices}
  \choice $x^{3}-x^{2}+C$
  \choice $3 x^{3}-x^{2}+3 x+C$
  \CorrectChoice $x^{3}-x^{2}+3 x+C$
  \choice $\dfrac{1}{2}\left(3 x^{2}-2 x+3\right)^{2}+C$
  \choice none of these
\end{choices}
\begin{solution}
Use, first, formula (2), then (3), replacing $u$ by $x$.
\end{solution}
 
\question
\[
\displaystyle\int\left(x-\dfrac{1}{2 x}\right)^{2} \,dx=
\]
\begin{choices}
  \choice $\dfrac{1}{3}\left(x-\dfrac{1}{2 x}\right)^{3}+C$
  \choice $x^{2}-1+\dfrac{1}{4 x^{2}}+C$
  \choice $\dfrac{x^{3}}{3}-2 x-\dfrac{1}{4 x}+C$
  \choice $\dfrac{x^{3}}{3}-x-\dfrac{4}{x}+C$
  \CorrectChoice none of these
\end{choices}
\begin{solution}
Hint: Expand. $\displaystyle\int\left(x^{2}-1+\dfrac{1}{4 x^{2}}\right) d x=\dfrac{x^{3}}{3}-x-\dfrac{1}{4 x}+C$.
\end{solution}

\question
\[
\displaystyle\int \sqrt{4-2 t} \,dt=
\]
\begin{choices}
  \CorrectChoice $-\dfrac{1}{3}(4-2 t)^{3 / 2}+C$
  \choice $\dfrac{2}{3}(4-2 t)^{3 / 2}+C$
  \choice $-\dfrac{1}{6}(4-2 t)^{3}+C$
  \choice $+\dfrac{1}{2}(4-2 t)^{2}+C$
  \choice $\dfrac{4}{3}(4-2 t)^{3 / 2}+C$
\end{choices}
\begin{solution}
By formula (3), with $u=4-2 t$ and $n=\dfrac{1}{2}$,
\[
\displaystyle\int \sqrt{4-2 t} d x=-\dfrac{1}{2} \displaystyle\int \sqrt{4-2 t} \cdot(-2 d t)=-\dfrac{1}{2} \dfrac{(4-2 t)^{3 / 2}}{3 / 2}+C
\]
\end{solution}





\newpage




\question
\[
\displaystyle\int(2-3 x)^{5} \,dx=
\]
\begin{choices}
  \choice $\dfrac{1}{6}(2-3 x)^{6}+C$
  \choice $-\dfrac{1}{2}(2-3 x)^{6}+C$
  \choice $\dfrac{1}{2}(2-3 x)^{6}+C$
  \CorrectChoice $-\dfrac{1}{18}(2-3 x)^{6}+C$
  \choice none of these
\end{choices}
\begin{solution}
Rewrite: $-\dfrac{1}{3} \displaystyle\int(2-3 x)^{5}(-3 d x)$
\end{solution}

\question
\[
\displaystyle\int \dfrac{1-3 y}{\sqrt{2 y-3 y^{2}}} \,dy=
\]
\begin{choices}
  \choice $4 \sqrt{2 y-3 y^{2}}+C$
  \choice $\dfrac{1}{4}\left(2 y-3 y^{2}\right)^{2}+C$
  \choice $\dfrac{1}{2} \ln \sqrt{2 y-3 y^{2}}+C$
  \choice $\dfrac{1}{4}\left(2 y-3 y^{2}\right)^{1 / 2}+C$
  \CorrectChoice $\sqrt{2 y-3 y^{2}}+C$
\end{choices}
\begin{solution}
Rewrite:
\[
\displaystyle\int\left(2 y-3 y^{2}\right)^{-1 / 2}(1-3 y) d y=\dfrac{1}{2} \displaystyle\int\left(2 y-3 y^{2}\right)^{-1 / 2}(2-6 y) d y
\]
Use (3).
\end{solution}






\newpage




\question
\[
\displaystyle\int \dfrac{d x}{3(2 x-1)^{2}}=
\]
\begin{choices}
  \choice $\dfrac{-3}{2 x-1}+C$
  \CorrectChoice $\dfrac{1}{6-12 x}+C$
  \choice $+\dfrac{6}{2 x-1}+C$
  \choice $\dfrac{2}{3 \sqrt{2 x-1}}+C$
  \choice $\dfrac{1}{3} \ln |2 x-1|+C$
\end{choices}
\begin{solution}
(B) Rewrite:

\[
\dfrac{1}{3} \displaystyle\int(2 x-1)^{-2} d x=\dfrac{1}{3} \cdot \dfrac{1}{2} \displaystyle\int(2 x-1)^{-2} \cdot 2 d x
\]

Using (3) yields $-\dfrac{1}{6(2 x-1)}+C$.
\end{solution}

\question
\[
\displaystyle\int \dfrac{2 \,du}{1+3 u}=
\]
\begin{choices}
  \CorrectChoice $\dfrac{2}{3} \ln |1+3 u|+C$
  \choice $-\dfrac{1}{3(1+3 u)^{2}}+C$
  \choice $2 \ln |1+3 u|+C$
  \choice $\dfrac{3}{(1+3 u)^{2}}+C$
  \choice none of these
\end{choices}
\begin{solution}
(A) This is equivalent to $\dfrac{1}{3} \cdot 2 \displaystyle\int \dfrac{3 d u}{1+3 u}$. Use (4).
\end{solution}






\newpage




\question
\[
\displaystyle\int \dfrac{t}{\sqrt{2 t^{2}-1}} \,dt=
\]
\begin{choices}
  \choice $\dfrac{1}{2} \ln \sqrt{2 t^{2}-1}+C$
  \choice $4 \ln \sqrt{2 t^{2}-1}+C$
  \choice $8 \sqrt{2 t^{2}-1}+C$
  \choice $-\dfrac{1}{4\left(2 t^{2}-1\right)}+C$
  \CorrectChoice $\dfrac{1}{2} \sqrt{2 t^{2}-1}+C$
\end{choices}
\begin{solution}
(E) Rewrite as $\dfrac{1}{4} \displaystyle\int\left(2 t^{2}-1\right)^{-1 / 2} \cdot 4 t d t$. Use (3).
\end{solution}

\question
\[
\displaystyle\int \cos 3 x \,dx=
\]
\begin{choices}
  \choice $3 \sin 3 x+C$
  \choice $-\sin 3 x+C$
  \choice $-\dfrac{1}{3} \sin 3 x+C$
  \CorrectChoice $\dfrac{1}{3} \sin 3 x+C$
  \choice $\dfrac{1}{2} \cos ^{2} 3 x+C$
\end{choices}
\begin{solution}
(D) Use (5) with $u=3 x ; d u=3 d x: \dfrac{1}{3} \displaystyle\int \cos (3 x)(3 d x)$
\end{solution}





\newpage





\question\label{calc05:4x2-first}
\[
\displaystyle\int \dfrac{x \,dx}{1+4 x^{2}}=
\]
\begin{choices}
  \CorrectChoice $\dfrac{1}{8} \ln \left(1+4 x^{2}\right)+C$
  \choice $\dfrac{1}{8\left(1+4 x^{2}\right)^{2}}+C$
  \choice $\dfrac{1}{4} \sqrt{1+4 x^{2}}+C$
  \choice $\dfrac{1}{2} \ln \left|1+4 x^{2}\right|+C$
  \choice $\dfrac{1}{2} \tan ^{-1} 2 x+C$
\end{choices}
\begin{solution}
(A) Use (4). If $u=1+4 x^{2}, d u=8 x d x: \dfrac{1}{8} \displaystyle\int \dfrac{8 x d x}{1+4 x^{2}}$
\end{solution}

\question
\[
\displaystyle\int \dfrac{d x}{1+4 x^{2}}=
\]
\begin{choices}
  \choice $\tan ^{-1}(2 x)+C$
  \choice $\dfrac{1}{8} \ln \left(1+4 x^{2}\right)+C$
  \choice $\dfrac{1}{8\left(1+4 x^{2}\right)^{2}}+C$
  \CorrectChoice $\dfrac{1}{2} \tan ^{-1}(2 x)+C$
  \choice $\dfrac{1}{8 x} \ln \left|1+4 x^{2}\right|+C$
\end{choices}
\begin{solution}
(D) Use (18). Let $u=2 x$; then $d u=2 d x: \dfrac{1}{2} \displaystyle\int \dfrac{2 d x}{1+(2 x)^{2}}$
\end{solution}




\newpage




\question
\[
\displaystyle\int \dfrac{x}{\left(1+4 x^{2}\right)^{2}} \,dx=
\]
\begin{choices}
  \choice $\dfrac{1}{8} \ln \left(1+4 x^{2}\right)^{2}+C$
  \choice $\dfrac{1}{4} \sqrt{1+4 x^{2}}+C$
  \CorrectChoice $-\dfrac{1}{8\left(1+4 x^{2}\right)}+C$
  \choice $-\dfrac{1}{3\left(1+4 x^{2}\right)^{3}}+C$
  \choice $-\dfrac{1}{\left(1+4 x^{2}\right)}+C$
\end{choices}
\begin{solution}
(C) Rewrite as $\dfrac{1}{8} \displaystyle\int\left(1+4 x^{3}\right)^{-2} \cdot(8 x d x)$. Use (3) with $n=-2$.
\end{solution}

\question\label{calc05:4x2-last}
\[
\displaystyle\int \dfrac{x \,dx}{\sqrt{1+4 x^{2}}}=
\]
\begin{choices}
  \choice $\dfrac{1}{8} \sqrt{1+4 x^{2}}+C$
  \CorrectChoice $\dfrac{\sqrt{1+4 x^{2}}}{4}+C$
  \choice $\dfrac{1}{2} \sin ^{-1} 2 x+C$
  \choice $\dfrac{1}{2} \tan ^{-1} 2 x+C$
  \choice $\dfrac{1}{8} \ln \sqrt{1+4 x^{2}}+C$
\end{choices}
\begin{solution}
(B) Rewrite as $\dfrac{1}{8} \displaystyle\int\left(1+4 x^{2}\right)^{-1 / 2} \cdot(8 x d x)$. Use (3) with $n=-\dfrac{1}{2}$.

\textit{Note carefully the differences in the integrands in Questions~\ref{calc05:4x2-first}--\ref{calc05:4x2-last}.}
\end{solution}





\newpage


\question\label{calc05:sqrt4y-first}
\[
\displaystyle\int \dfrac{d y}{\sqrt{4-y^{2}}}=
\]
\begin{choices}
  \choice $\dfrac{1}{2} \sin ^{-1} \dfrac{y}{2}+C$
  \choice $-\sqrt{4-y^{2}}+C$
  \CorrectChoice $\sin ^{-1} \dfrac{y}{2}+C$
  \choice $-\dfrac{1}{2} \ln \sqrt{4-y^{2}}+C$
  \choice $-\dfrac{1}{3\left(4-y^{2}\right)^{3 / 2}}+C$
\end{choices}
\begin{solution}
(C) Use (17); rewrite as $\displaystyle\int \dfrac{\dfrac{1}{2} d y}{\sqrt{1-\left(\dfrac{y}{2}\right)^{2}}}$.
\end{solution}

\question\label{calc05:sqrt4y-last}
\[
\displaystyle\int \dfrac{y \,dy}{\sqrt{4-y^{2}}}=
\]
\begin{choices}
  \choice $\dfrac{1}{2} \sin ^{-1} \dfrac{y}{2}+C$
  \CorrectChoice $-\sqrt{4-y^{2}}+C$
  \choice $\sin ^{-1} \dfrac{y}{2}+C$
  \choice $-\dfrac{1}{2} \ln \sqrt{4-y^{2}}+C$
  \choice $2 \sqrt{4-y^{2}}+C$
\end{choices}
\begin{solution}
(B) Rewrite as $-\dfrac{1}{2} \displaystyle\int\left(4-y^{2}\right)^{-1 / 2} \cdot(-2 y d y)$. Use (3).

\textit{Compare the integrands in Questions~\ref{calc05:sqrt4y-first} and~\ref{calc05:sqrt4y-last}, noting the difference.}
\end{solution}


\newpage





\question
\[
\displaystyle\int \dfrac{2 x+1}{2 x} \,dx=
\]
\begin{choices}
  \CorrectChoice $x+\dfrac{1}{2} \ln |x|+C$
  \choice $1+\dfrac{1}{2} x^{-1}+C$
  \choice $x+2 \ln |x|+C$
  \choice $x+\ln |2 x|+C$
  \choice $\dfrac{1}{2}\left(2 x-\dfrac{1}{x^{2}}\right)+C$
\end{choices}
\begin{solution}
(A) Divide to obtain $\displaystyle\int\left(1+\dfrac{1}{2} \cdot \dfrac{1}{x}\right) d x$. Use (2), (3), and (4). Remember that $\displaystyle\int k d x=k x+C$ whenever $k \neq 0$.
\end{solution}

\question
\[
\displaystyle\int \dfrac{(x-2)^{3}}{x^{2}} \,dx=
\]
\begin{choices}
  \choice $\dfrac{(x-2)^{4}}{4 x^{2}}+C$
  \choice $\dfrac{x^{2}}{2}-6 x+6 \ln |x|-\dfrac{8}{x}+C$
  \choice $\dfrac{x^{2}}{2}-3 x+6 \ln |x|+\dfrac{4}{x}+C$
  \choice $-\dfrac{(x-2)^{4}}{4 x}+C$
  \CorrectChoice none of these
\end{choices}
\begin{solution}
(E) $\quad \displaystyle\int \dfrac{(x-2)^{3}}{x^{2}} d x=\displaystyle\int\left(x-6+\dfrac{12}{x}-\dfrac{8}{x^{2}}\right) d x=\dfrac{x^{2}}{2}-6 x+12 \ln |x|+\dfrac{8}{x}+C$. We used the Binomial Theorem on page 667 with $n=3$ to expand $(x-2)^{3}$.
\end{solution}





\newpage


\question
\[
\displaystyle\int\left(\sqrt{t}-\dfrac{1}{\sqrt{t}}\right)^{2} \,dt=
\]
\begin{choices}
  \choice $t-2+\dfrac{1}{t}+C$
  \choice $\dfrac{t^{3}}{3}-2 t-\dfrac{1}{t}+C$
  \choice $\dfrac{t^{2}}{2}+\ln |t|+C$
  \CorrectChoice $\dfrac{t^{2}}{2}-2 t+\ln |t|+C$
  \choice $\dfrac{t^{2}}{2}-t-\dfrac{1}{t^{2}}+C$
\end{choices}
\begin{solution}
(D) The integral is equivalent to $\displaystyle\int\left(t-2+\dfrac{1}{t}\right) d t$. Integrate term by term.
\end{solution}

\question
\[
\displaystyle\int\left(4 x^{1 / 3}-5 x^{3 / 2}-x^{-1 / 2}\right) \,dx=
\]
\begin{choices}
  \CorrectChoice $3 x^{4 / 3}-2 x^{5 / 2}-2 x^{1 / 2}+C$
  \choice $3 x^{4 / 3}-2 x^{5 / 2}+2 x^{1 / 2}+C$
  \choice $6 x^{2 / 3}-2 x^{5 / 2}-\dfrac{1}{2} x^{2}+C$
  \choice $\dfrac{4}{3} x^{-2 / 3}-\dfrac{15}{2} x^{1 / 2}+\dfrac{1}{2} x^{-3 / 2}+C$
  \choice none of these
\end{choices}
\begin{solution}
(A) Integrate term by term.
\end{solution}




\newpage


\question
\[
\displaystyle\int \dfrac{x^{3}-x-1}{x^{2}} \,dx=
\]
\begin{choices}
  \choice $\dfrac{\dfrac{1}{4} x^{4}-\dfrac{1}{2} x^{2}-x}{\dfrac{1}{3} x^{3}}+C$
  \choice $1+\dfrac{1}{x^{2}}+\dfrac{2}{x^{3}}+C$
  \choice $\dfrac{x^{2}}{2}-\ln |x|-\dfrac{1}{x}+C$
  \CorrectChoice $\dfrac{x^{2}}{2}-\ln |x|+\dfrac{1}{x}+C$
  \choice $\dfrac{x^{2}}{2}-\ln |x|+\dfrac{2}{x^{3}}+C$
\end{choices}
\begin{solution}
(D) Division yields

\[
\displaystyle\int\left(x-\dfrac{1}{x}-\dfrac{1}{x^{2}}\right) d x=\displaystyle\int x d x-\displaystyle\int \dfrac{1}{x} d x-\displaystyle\int \dfrac{1}{x^{2}} d x
\]
\end{solution}

\question
\[
\displaystyle\int \dfrac{d y}{\sqrt{y}(1-\sqrt{y})}=
\]
\begin{choices}
  \choice $4 \sqrt{1-\sqrt{y}}+C$
  \choice $\dfrac{1}{2} \ln |1-\sqrt{y}|+C$
  \choice $2 \ln (1-\sqrt{y})+C$
  \choice $2 \sqrt{y}-\ln |y|+C$
  \CorrectChoice $-2 \ln |1-\sqrt{y}|+C$
\end{choices}
\begin{solution}
Use formula (4) with $u=1-\sqrt{y}=1-y^{1 / 2}$. Then $d u=-\dfrac{1}{2 \sqrt{y}} d y$. Note that the integral can be written as $-2 \displaystyle\int \dfrac{1}{(1-\sqrt{y})}\left(-\dfrac{1}{2 \sqrt{y}}\right) d y$.
\end{solution}




\newpage


\question
\[
\displaystyle\int \dfrac{u \,du}{\sqrt{4-9 u^{2}}}=
\]
\begin{choices}
  \choice $\dfrac{1}{3} \sin ^{-1} \dfrac{3 u}{2}+C$
  \choice $-\dfrac{1}{18} \ln \sqrt{4-9 u^{2}}+C$
  \choice $2 \sqrt{4-9 u^{2}}+C$
  \choice $\dfrac{1}{6} \sin ^{-1} \dfrac{3}{2} u+C$
  \CorrectChoice $-\dfrac{1}{9} \sqrt{4-9 u^{2}}+C$
\end{choices}
\begin{solution}
Rewrite as $-\dfrac{1}{18} \displaystyle\int\left(4-9 u^{2}\right)^{-1 / 2}(-18 u d u)$ and use formula (3).
\end{solution}

\question\label{calc05:sin-theta-cos}
\[
\displaystyle\int \sin \theta \cos \theta d \theta=
\]
\begin{choices}
  \choice $-\dfrac{\sin ^{2} \theta}{2}+C$
  \CorrectChoice $-\dfrac{1}{4} \cos 2 \theta+C$
  \choice $\dfrac{\cos ^{2} \theta}{2}+C$
  \choice $\dfrac{1}{2} \sin 2 \theta+C$
  \choice $\cos 2 \theta+C$
\end{choices}
\begin{solution}
The integral is equal to $\dfrac{1}{2} \displaystyle\int \sin 2 \theta d \theta$. Use formula (6) with $u=2 \theta ; d u=2 d \theta$.
\end{solution}

\question
\[
\displaystyle\int \dfrac{\sin \sqrt{x}}{\sqrt{x}} \,dx=
\]
\begin{choices}
  \choice $-2 \cos ^{1 / 2} x+C$
  \choice $-\cos \sqrt{x}+C$
  \CorrectChoice $-2 \cos \sqrt{x}+C$
  \choice $\dfrac{3}{2} \sin ^{3 / 2} x+C$
  \choice $\dfrac{1}{2} \cos \sqrt{x}+C$
\end{choices}
\begin{solution}
Use formula (6) with $u=\sqrt{x} ; d u=\dfrac{1}{2 \sqrt{x}} d x ; 2 \displaystyle\int \sin (\sqrt{x})\left(\dfrac{1}{2 \sqrt{x}} d x\right)$
\end{solution}




\newpage


\question
\[
\displaystyle\int t \cos (2 t)^{2} \,dt=
\]
\begin{choices}
  \CorrectChoice $\dfrac{1}{8} \sin \left(4 t^{2}\right)+C$
  \choice $\dfrac{1}{2} \cos ^{2}(2 t)+C$
  \choice $-\dfrac{1}{8} \sin \left(4 t^{2}\right)+C$
  \choice $\dfrac{1}{4} \sin (2 t)^{2}+C$
  \choice none of these
\end{choices}
\begin{solution}
Use formula (5) with $u=4 t^{2} ; d u=8 t d t ; \dfrac{1}{8} \displaystyle\int \cos \left(4 t^{2}\right)(8 t d t)$
\end{solution}

\question
\[
\displaystyle\int \cos ^{2} 2 x \,dx=
\]
\begin{choices}
  \CorrectChoice $\dfrac{x}{2}+\dfrac{\sin 4 x}{8}+C$
  \choice $\dfrac{x}{2}-\dfrac{\sin 4 x}{8}+C$
  \choice $\dfrac{x}{4}+\dfrac{\sin 4 x}{4}+C$
  \choice $\dfrac{x}{4}+\dfrac{\sin 4 x}{16}+C$
  \choice $\dfrac{1}{4}(x+\sin 4 x)+C$
\end{choices}
\begin{solution}
Using the Half-Angle Formula (23) on page 669 with $\alpha=2 x$ yields $\displaystyle\int\left(\dfrac{1}{2}+\dfrac{1}{2} \cos 4 x\right) d x$.
\end{solution}

\question
\[
\displaystyle\int \sin 2 \theta d \theta=
\]
\begin{choices}
  \choice $\dfrac{1}{2} \cos 2 \theta+C$
  \choice $-2 \cos 2 \theta+C$
  \choice $-\sin ^{2} \theta+C$
  \choice $\cos ^{2} \theta+C$
  \CorrectChoice $-\dfrac{1}{2} \cos 2 \theta+C$
\end{choices}
\begin{solution}
Use formula (6): $\dfrac{1}{2} \displaystyle\int \sin 2 \theta(2 d \theta)$. (See Question~\ref{calc05:sin-theta-cos}.)
\end{solution}




\newpage


\question
\[
\displaystyle\int x \cos x \,dx=
\]
\begin{choices}
  \choice $x \sin x+C$
  \CorrectChoice $x \sin x+\cos x+C$
  \choice $x \sin x-\cos x+C$
  \choice $\cos x-x \sin x+C$
  \choice $\dfrac{x^{2}}{2} \sin x+C$
\end{choices}
\begin{solution}
Integrate by parts (page 224). Let $u=x$ and $d v=\cos x d x$. Then $d u=d x$ and $v=\sin x$. The given integral equals $x \sin x-\displaystyle\int \sin x d x$.
\end{solution}

\question
\[
\displaystyle\int \dfrac{d u}{\cos ^{2} 3 u}=
\]
\begin{choices}
  \choice $-\dfrac{\sec 3 u}{3}+C$
  \choice $\tan 3 u+C$
  \choice $u+\dfrac{\sec 3 u}{3}+C$
  \CorrectChoice $\dfrac{1}{3} \tan 3 u+C$
  \choice $\dfrac{1}{3 \cos 3 u}+C$
\end{choices}
\begin{solution}
Replace $\dfrac{1}{\cos ^{2} 3 u}$ by $\sec ^{2} 3 u$; then use formula (9): $\dfrac{1}{3} \displaystyle\int \sec ^{2} 3 u(3 d u)$
\end{solution}

\question
\[
\displaystyle\int \dfrac{\cos x \,dx}{\sqrt{1+\sin x}}=
\]
\begin{choices}
  \choice $-\dfrac{1}{2}(1+\sin x)^{1 / 2}+C$
  \choice $\ln \sqrt{1+\sin x}+C$
  \CorrectChoice $2 \sqrt{1+\sin x}+C$
  \choice $\ln |1+\sin x|+C$
  \choice $\dfrac{2}{3(1+\sin x)^{3 / 2}}+C$
\end{choices}
\begin{solution}
Rewrite using $u=1+\sin x$ and $d u=\cos x d x$ as $\displaystyle\int(1+\sin x)^{-1 / 2}(\cos x d x)$.
Use formula (3).
\end{solution}




\newpage


\question
\[
\displaystyle\int \dfrac{\cos (\theta-1) d \theta}{\sin ^{2}(\theta-1)}=
\]
\begin{choices}
  \choice $2 \ln \sin |\theta-1|+C$
  \CorrectChoice $-\csc (\theta-1)+C$
  \choice $-\dfrac{1}{3} \sin ^{-3}(\theta-1)+C$
  \choice $-\cot (\theta-1)+C$
  \choice $\csc (\theta-1)+C$
\end{choices}
\begin{solution}
(B) The integral is equivalent to $\displaystyle\int \csc (\theta-1) \cot (\theta-1) d \theta$. Use formula (12).
\end{solution}

\question
\[
\displaystyle\int \sec \dfrac{t}{2} \,dt=
\]
\begin{choices}
  \choice $\ln \left|\sec \dfrac{t}{2}+\tan \dfrac{t}{2}\right|+C$
  \choice $2 \tan ^{2} \dfrac{t}{2}+C$
  \choice $2 \ln \cos \dfrac{t}{2}+C$
  \choice $\ln |\sec t+\tan t|+C$
  \CorrectChoice $2 \ln \left|\sec \dfrac{t}{2}+\tan \dfrac{t}{2}\right|+C$
\end{choices}
\begin{solution}
(E) Use formula (13) with $u=\dfrac{t}{2} ; d u=\dfrac{1}{2} d t: 2 \displaystyle\int \sec \dfrac{t}{2}\left(\dfrac{1}{2} d t\right)$
\end{solution}

\question
\[
\displaystyle\int \dfrac{\sin 2 x \,dx}{\sqrt{1+\cos ^{2} x}}=
\]
\begin{choices}
  \CorrectChoice $-2 \sqrt{1+\cos ^{2} x}+C$
  \choice $\dfrac{1}{2} \ln \left(1+\cos ^{2} x\right)+C$
  \choice $\sqrt{1+\cos ^{2} x}+C$
  \choice $-\ln \sqrt{1+\cos ^{2} x}+C$
  \choice $2 \ln |\sin x|+C$
\end{choices}
\begin{solution}
(A) Replace $\sin 2 x$ by $2 \sin x \cos x$; then the integral is equivalent to
\[
-\displaystyle\int \dfrac{-2 \sin x \cos x}{\sqrt{1+\cos ^{2} x}} d x=-\displaystyle\int u^{-1 / 2} d u
\]
where $u=1+\cos ^{2} x$ and $d u=-2 \sin x \cos x d x$. Use formula (3).
\end{solution}




\newpage


\question
\[
\displaystyle\int \sec ^{3 / 2} x \tan x \,dx=
\]
\begin{choices}
  \choice $\dfrac{2}{5} \sec ^{5 / 2} x+C$
  \choice $-\dfrac{2}{3} \cos ^{-3 / 2} x+C$
  \choice $\sec ^{3 / 2} x+C$
  \CorrectChoice $\dfrac{2}{3} \sec ^{3 / 2} x+C$
  \choice none of these
\end{choices}
\begin{solution}
(D) Rewriting in terms of sines and cosines yields

\[
\displaystyle\int \dfrac{\sin x}{\cos ^{5 / 2} x} d x=-\displaystyle\int \cos ^{-5 / 2} x(-\sin x) d x=-\left(-\dfrac{2}{3}\right) \cos ^{-3 / 2} x+C
\]
\end{solution}

\question
\[
\displaystyle\int \tan \theta d \theta=
\]
\begin{choices}
  \choice $-\ln |\sec \theta|+C$
  \choice $\sec ^{2} \theta+C$
  \choice $\ln |\sin \theta|+C$
  \choice $\sec \theta+C$
  \CorrectChoice $-\ln |\cos \theta|+C$
\end{choices}
\begin{solution}
Use formula (7).
\end{solution}

\question
\[
\displaystyle\int \dfrac{d x}{\sin ^{2} 2 x}=
\]
\begin{choices}
  \choice $\dfrac{1}{2} \csc 2 x \cot 2 x+C$
  \choice $-\dfrac{2}{\sin 2 x}+C$
  \CorrectChoice $-\dfrac{1}{2} \cot 2 x+C$
  \choice $-\cot x+C$
  \choice $-\csc 2 x+C$
\end{choices}
\begin{solution}
Replace $\dfrac{1}{\sin ^{2} 2 x}$ by $\csc ^{2} 2 x$ and use formula (10): $\dfrac{1}{2} \displaystyle\int \csc ^{2} 2 x(2 d x)$
\end{solution}




\newpage


\question
\[
\displaystyle\int \dfrac{\tan ^{-1} y}{1+y^{2}} \,dy=
\]
\begin{choices}
  \choice $\sec ^{-1} y+C$
  \choice $\left(\tan ^{-1} y\right)^{2}+C$
  \choice $\ln \left(1+y^{2}\right)+C$
  \choice $\ln \left(\tan ^{-1} y\right)+C$
  \CorrectChoice none of these
\end{choices}
\begin{solution}
Let $u=\tan ^{-1} y$; then integrate $\displaystyle\int u d u$. The correct answer is $\dfrac{1}{2}\left(\tan ^{-1} y\right)^{2}+C$.
\end{solution}

\question
\[
\displaystyle\int \sin 2 \theta \cos \theta d \theta=
\]
\begin{choices}
  \CorrectChoice $-\dfrac{2}{3} \cos ^{3} \theta+C$
  \choice $\dfrac{2}{3} \cos ^{3} \theta+C$
  \choice $\sin ^{2} \theta \cos \theta+C$
  \choice $\cos ^{3} \theta+C$
  \choice none of these
\end{choices}
\begin{solution}
Replacing $\sin 2 \theta$ by $2 \sin \theta \cos \theta$ yields


\[
\displaystyle\int 2 \sin \theta \cos ^{2} \theta d \theta=-2 \displaystyle\int(\cos \theta)^{2}(-\sin \theta d \theta)=-\dfrac{2}{3} \cos ^{3} \theta+C
\]
\end{solution}

\question
\[
\displaystyle\int \dfrac{\sin 2 t}{1-\cos 2 t} \,dt=
\]
\begin{choices}
  \choice $\dfrac{2}{(1-\cos 2 t)^{2}}+C$
  \choice $-\ln |1-\cos 2 t|+C$
  \CorrectChoice $\ln \sqrt{|1-\cos 2 t|}+C$
  \choice $\sqrt{1-\cos 2 t}+C$
  \choice $2 \ln |1-\cos 2 t|+C$
\end{choices}
\begin{solution}
$\dfrac{1}{2} \displaystyle\int \dfrac{2 \sin 2 t d t}{1-\cos 2 t}=\dfrac{1}{2} \ln |1-\cos 2 t|+C$.
\end{solution}


\newpage


\question
\[
\displaystyle\int \cot 2 u \,du=
\]
\begin{choices}
  \choice $\ln |\sin u|+C$
  \CorrectChoice $\dfrac{1}{2} \ln |\sin 2 u|+C$
  \choice $-\dfrac{1}{2} \csc ^{2} 2 u+C$
  \choice $-\sec 2 u+C$
  \choice $2 \ln |\sin 2 u|+C$
\end{choices}
\begin{solution}
Rewrite as $\dfrac{1}{2} \displaystyle\int \cot 2 u(2 d u)$ and use formula (8).
\end{solution}

\question
\[
\displaystyle\int \dfrac{e^{x}}{e^{x}-1} \,dx=
\]
\begin{choices}
  \choice $x+\ln \left|e^{x}-1\right|+C$
  \choice $x-e^{x}+C$
  \choice $x-\dfrac{1}{\left(e^{x}-1\right)^{2}}+C$
  \choice $1+\dfrac{1}{e^{x}-1}+C$
  \CorrectChoice $\ln \left|e^{x}-1\right|+C$
\end{choices}
\begin{solution}
Use formula (4) with $u=e^{x}-1 ; d u=e^{x} d x$.
\end{solution}



\newpage


\question
\[
\displaystyle\int \dfrac{x-1}{x(x-2)} \,dx=
\]
\begin{choices}
  \choice $\dfrac{1}{2} \ln |x|+\ln |x-2|+C$
  \choice $\dfrac{1}{2} \ln \left|\dfrac{x-2}{x}\right|+C$
  \choice $\ln |x-2|+\ln |x|+C$
  \CorrectChoice $\dfrac{1}{2} \ln |x(x-2)|+C$
  \choice none of these
\end{choices}
\begin{solution}
Use partial fractions; find $A$ and $B$ such that


\[
\dfrac{x-1}{x(x-2)}=\dfrac{A}{x}+\dfrac{B}{x-2}
\]

Then $x-1=A(x-2)+B x$.
Set $x=0$ : $-1=-2 A$ and $A=\dfrac{1}{2}$.\\
Set $x=2: 1=2 B$ and $B=\dfrac{1}{2}$.\\
So the given integral equals

\[
\begin{aligned}
\displaystyle\int\left(\dfrac{1}{2 x}+\dfrac{1}{2(x-2)}\right) d x & =\dfrac{1}{2} \ln |x|+\dfrac{1}{2} \ln |x-2|+C \\
& =\dfrac{1}{2} \ln |x(x-2)|+C
\end{aligned}
\]
\end{solution}

\question
\[
\displaystyle\int x e^{x^{2}} \,dx=
\]
\begin{choices}
  \CorrectChoice $\dfrac{1}{2} e^{x^{2}}+C$
  \choice $e^{x^{2}}\left(2 x^{2}+1\right)+C$
  \choice $2 e^{x^{2}}+C$
  \choice $e^{x^{2}}+C$
  \choice $\dfrac{1}{2} e^{x^{2}+1}+C$
\end{choices}
\begin{solution}
Use formula (15) with $u=x^{2} ; d u=2 x d x ; \dfrac{1}{2} \displaystyle\int e^{x^{2}}(2 x d x)$.
\end{solution}



\newpage


\question
\[
\displaystyle\int \cos \theta e^{\sin \theta} d \theta=
\]
\begin{choices}
  \choice $e^{\sin \theta+1}+C$
  \CorrectChoice $e^{\sin \theta}+C$
  \choice $-e^{\sin \theta}+C$
  \choice $e^{\cos \theta}+C$
  \choice $e^{\sin \theta}(\cos \theta-\sin \theta)+C$
\end{choices}
\begin{solution}
Use formula (15) with $u=\sin \theta ; d u=\cos \theta d \theta$.
\end{solution}

\question
\[
\displaystyle\int e^{2 \theta} \sin e^{2 \theta} d \theta=
\]
\begin{choices}
  \choice $\cos e^{2 \theta}+C$
  \choice $2 e^{4 \theta}\left(\cos e^{2 \theta}+\sin e^{2 \theta}\right)+C$
  \CorrectChoice $-\dfrac{1}{2} \cos e^{2 \theta}+C$
  \choice $-2 \cos e^{2 \theta}+C$
  \choice none of these
\end{choices}
\begin{solution}
Use formula (6) with $u=e^{2 \theta} ; d u=2 e^{2 \theta} d \theta ; \dfrac{1}{2} \displaystyle\int \sin e^{2 \theta}\left(2 e^{2 \theta} d \theta\right)$.
\end{solution}

\question
\[
\displaystyle\int \dfrac{e^{\sqrt{x}}}{\sqrt{x}} \,dx=
\]
\begin{choices}
  \choice $2 \sqrt{x}\left(e^{\sqrt{x}}-1\right)+C$
  \CorrectChoice $2 e^{\sqrt{x}}+C$
  \choice $\dfrac{e^{\sqrt{x}}}{2}\left(\dfrac{1}{x}+\dfrac{1}{x \sqrt{x}}\right)+C$
  \choice $\dfrac{1}{2} e^{\sqrt{x}}+C$
  \choice none of these
\end{choices}
\begin{solution}
Use formula (15) with $u=\sqrt{x}=x^{1 / 2} ; d u=\dfrac{1}{2 \sqrt{x}} d x$.
\end{solution}



\newpage


\question
\[
\displaystyle\int x e^{-x} \,dx=
\]
\begin{choices}
  \choice $e^{-x}(1-x)+C$
  \choice $\dfrac{e^{1-x}}{1-x}+C$
  \CorrectChoice $-e^{-x}(x+1)+C$
  \choice $-\dfrac{x^{2}}{2} e^{-x}+C$
  \choice $e^{-x}(x+1)+C$
\end{choices}
\begin{solution}
Use the Parts Formula. Let $u=x, d v=e^{-x} d x ; d u=d x, v=-e^{-x}$. Then,


\[
-x e^{-x}+\displaystyle\int e^{-x} d x=-x e^{-x}-e^{-x}+C
\]
\end{solution}

\question
\[
\displaystyle\int x^{2} e^{x} \,dx=
\]
\begin{choices}
  \choice $e^{x}\left(x^{2}+2 x\right)+C$
  \choice $e^{x}\left(x^{2}-2 x-2\right)+C$
  \CorrectChoice $e^{x}\left(x^{2}-2 x+2\right)+C$
  \choice $e^{x}(x-1)^{2}+C$
  \choice $e^{x}(x+1)^{2}+C$
\end{choices}
\begin{solution}
See Example 44, page 225.
\end{solution}

\question
\[
\displaystyle\int \dfrac{e^{x}+e^{-x}}{e^{x}-e^{-x}} \,dx=
\]
\begin{choices}
  \choice $x-\ln \left|e^{x}-e^{-x}\right|+C$
  \choice $x+2 \ln \left|e^{x}-e^{-x}\right|+C$
  \choice $-\dfrac{1}{2}\left(e^{x}-e^{-x}\right)^{-2}+C$
  \CorrectChoice $\ln \left|e^{x}-e^{-x}\right|+C$
  \choice $\ln \left(e^{x}+e^{-x}\right)+C$
\end{choices}
\begin{solution}
The integral is of the form $\displaystyle\int \dfrac{d u}{u}$; use (4).
\end{solution}



\newpage


\question
\[
\displaystyle\int \dfrac{e^{x}}{1+e^{2 x}} \,dx=
\]
\begin{choices}
  \CorrectChoice $\tan ^{-1} e^{x}+C$
  \choice $\dfrac{1}{2} \ln \left(1+e^{2 x}\right)+C$
  \choice $\ln \left(1+e^{2 x}\right)+C$
  \choice $\dfrac{1}{2} \tan ^{-1} e^{x}+C$
  \choice $2 \tan ^{-1} e^{x}+C$
\end{choices}
\begin{solution}
The integral has the form $\displaystyle\int \dfrac{d u}{1+u^{2}}$. Use formula (18), with $u=e^{x}, d u= e^{x} d x$.
\end{solution}

\question
\[
\displaystyle\int \dfrac{\ln v d v}{v}=
\]
\begin{choices}
  \choice $\ln |\ln v|+C$
  \choice $\ln \dfrac{v^{2}}{2}+C$
  \CorrectChoice $\dfrac{1}{2}(\ln v)^{2}+C$
  \choice $2 \ln v+C$
  \choice $\dfrac{1}{2} \ln v^{2}+C$
\end{choices}
\begin{solution}
Let $u=\ln v$; then $d u=\dfrac{d v}{v}$. Use formula (3) for $\displaystyle\int \ln v\left(\dfrac{1}{v} d v\right)$.
\end{solution}

\question
\[
\displaystyle\int \dfrac{\ln \sqrt{x}}{x} \,dx=
\]
\begin{choices}
  \choice $\dfrac{\ln ^{2} \sqrt{x}}{\sqrt{x}}+C$
  \choice $\ln ^{2} x+C$
  \choice $\dfrac{1}{2} \ln |\ln x|+C$
  \choice $\dfrac{(\ln \sqrt{x})^{2}}{2}+C$
  \CorrectChoice $\dfrac{1}{4} \ln ^{2} x+C$
\end{choices}
\begin{solution}
Hint: $\ln \sqrt{x}=\dfrac{1}{2} \ln x$; the integral is $\dfrac{1}{2} \displaystyle\int(\ln x)\left(\dfrac{1}{x} d x\right)$
\end{solution}



\newpage


\question
\[
\displaystyle\int x^{3} \ln x \,dx=
\]
\begin{choices}
  \choice $x^{2}(3 \ln x+1)+C$
  \CorrectChoice $\dfrac{x^{4}}{16}(4 \ln x-1)+C$
  \choice $\dfrac{x^{4}}{4}(\ln x-1)+C$
  \choice $3 x^{2}\left(\ln x-\dfrac{1}{2}\right)+C$
  \choice none of these
\end{choices}
\begin{solution}
Use parts, letting $u=\ln x$ and $d v=x^{3} d x$. Then $d u=\dfrac{1}{x} d x$ and $v=\dfrac{x^{4}}{4}$. The integral equals $\dfrac{x^{4}}{4} \ln x-\dfrac{1}{4} \displaystyle\int x^{3} d x$.
\end{solution}

\question\label{calc05:ln-eta}
\[
\displaystyle\int \ln \eta d \eta=
\]
\begin{choices}
  \choice $\dfrac{1}{2} \ln ^{2} \eta+C$
  \CorrectChoice $\eta(\ln \eta-1)+C$
  \choice $\dfrac{1}{2} \ln \eta^{2}+C$
  \choice $\ln \eta(\eta-1)+C$
  \choice $\eta \ln \eta+\eta+C$
\end{choices}
\begin{solution}
Use parts, letting $u=\ln \eta$ and $d v=d x$. Then $d u=\dfrac{1}{\eta} d \eta$ and $v=\eta$. The integral equals $\eta \ln \eta-\displaystyle\int d \eta$.
\end{solution}

\question
\[
\displaystyle\int \ln x^{3} \,dx=
\]
\begin{choices}
  \choice $\dfrac{3}{2} \ln ^{2} x+C$
  \CorrectChoice $3 x(\ln x-1)+C$
  \choice $3 \ln x(x-1)+C$
  \choice $\dfrac{3 x \ln ^{2} x}{2}+C$
  \choice none of these
\end{choices}
\begin{solution}
Rewrite $\ln x^{3}$ as $3 \ln x$, and use the method of Question~\ref{calc05:ln-eta}.
\end{solution}



\newpage


\question
\[
\displaystyle\int \dfrac{\ln y}{y^{2}} \,dy=
\]
\begin{choices}
  \choice $\dfrac{1}{y}(1-\ln y)+C$
  \choice $\dfrac{1}{2 y} \ln ^{2} y+C$
  \choice $-\dfrac{1}{3 y^{3}}(4 \ln y+1)+C$
  \CorrectChoice $-\dfrac{1}{y}(\ln y+1)+C$
  \choice $\dfrac{\ln y}{y}-\dfrac{1}{y}+C$
\end{choices}
\begin{solution}
Use parts, letting $u=\ln y$ and $d v=y^{-2} d y$. Then $d u=\dfrac{1}{y} d y$ and $v=-\dfrac{1}{y}$. The Parts Formula yields $\dfrac{-\ln y}{y}+\displaystyle\int \dfrac{1}{y^{2}} d y$.
\end{solution}

\question
\[
\displaystyle\int \dfrac{d v}{v \ln v}=
\]
\begin{choices}
  \choice $\dfrac{1}{\ln v^{2}}+C$
  \choice $-\dfrac{1}{\ln ^{2} v}+C$
  \choice $-\ln |\ln v|+C$
  \choice $\ln \dfrac{1}{v}+C$
  \CorrectChoice $\ln |\ln v|+C$
\end{choices}
\begin{solution}
The integral has the form $\displaystyle\int \dfrac{d u}{u}$, where $u=\ln v: \displaystyle\int \dfrac{\left(\dfrac{1}{v} d v\right)}{\ln v}$
\end{solution}



\newpage


\question
\[
\displaystyle\int \dfrac{y-1}{y+1} \,dy=
\]
\begin{choices}
  \CorrectChoice $y-2 \ln |y+1|+C$
  \choice $1-\dfrac{2}{y+1}+C$
  \choice $\ln \dfrac{|y|}{(y+1)^{2}}+C$
  \choice $1-2 \ln |y+1|+C$
  \choice $\ln \left|\dfrac{e^{y}}{y+1}\right|+C$
\end{choices}
\begin{solution}
By long division, the integrand is equivalent to $1-\dfrac{2}{y+1}$.
\end{solution}

\question
\[
\displaystyle\int \dfrac{d x}{x^{2}+2 x+2}=
\]
\begin{choices}
  \choice $\ln \left(x^{2}+2 x+2\right)+C$
  \choice $\ln |x+1|+C$
  \CorrectChoice $\arctan (x+1)+C$
  \choice $\dfrac{1}{\dfrac{1}{3} x^{3}+x^{2}+2 x}+C$
  \choice $-\dfrac{1}{x}+\dfrac{1}{2} \ln |x|+\dfrac{x}{2}+C$
\end{choices}
\begin{solution}
$\displaystyle\int \dfrac{d x}{(x+1)^{2}+1}=\displaystyle\int \dfrac{d x}{1+(x+1)^{2}}$; use formula (18) with $u=x+1$.
\end{solution}

\question
\[
\displaystyle\int \sqrt{x}(\sqrt{x}-1) \,dx=
\]
\begin{choices}
  \choice $2\left(x^{3 / 2}-x\right)+C$
  \choice $\dfrac{x^{2}}{2}-x+C$
  \choice $\dfrac{1}{2}(\sqrt{x}-1)^{2}+C$
  \CorrectChoice $\dfrac{1}{2} x^{2}-\dfrac{2}{3} x^{3 / 2}+C$
  \choice $x-2 \sqrt{x}+C$
\end{choices}
\begin{solution}
Multiply to get $\displaystyle\int(x-\sqrt{x}) d x$.
\end{solution}



\newpage



\question\label{calc05:bc-first}
\[
\displaystyle\int e^{\theta} \cos \theta d \theta=
\]
\begin{choices}
  \choice $e^{\theta}(\cos \theta-\sin \theta)+C$
  \choice $e^{\theta} \sin \theta+C$
  \CorrectChoice $\dfrac{1}{2} e^{\theta}(\sin \theta+\cos \theta)+C$
  \choice $2 e^{\theta}(\sin \theta+\cos \theta)+C$
  \choice $\dfrac{1}{2} e^{\theta}(\sin \theta-\cos \theta)+C$
\end{choices}
\begin{solution}
See Example 45, page 225. Replace $x$ by $\theta$.
\end{solution}

\question
\[
\displaystyle\int \dfrac{(1-\ln t)^{2}}{t} \,dt=
\]
\begin{choices}
  \choice $\dfrac{1}{3}(1-\ln t)^{3}+C$
  \choice $\ln t-2 \ln ^{2} t+\ln ^{3} t+C$
  \choice $-2(1-\ln t)+C$
  \choice $\ln t-\ln ^{2} t+\dfrac{\ln t^{3}}{3}+C$
  \CorrectChoice $-\dfrac{(1-\ln t)^{3}}{3}+C$
\end{choices}
\begin{solution}
The integral equals $-\displaystyle\int(1-\ln t)^{2}\left(-\dfrac{1}{t} d t\right)$; it is equivalent to $-\displaystyle\int u^{2} d u$, where $u=1-\ln t$.
\end{solution}

\question
\[
\displaystyle\int u \sec ^{2} u \,du=
\]
\begin{choices}
  \CorrectChoice $u \tan u+\ln |\cos u|+C$
  \choice $\dfrac{u^{2}}{2} \tan u+C$
  \choice $\dfrac{1}{2} \sec u \tan u+C$
  \choice $u \tan u-\ln |\sin u|+C$
  \choice $u \sec u-\ln |\sec u+\tan u|+C$
\end{choices}
\begin{solution}
Replace $u$ by $x$ in the given integral to avoid confusion in applying the Parts Formula. To integrate $\displaystyle\int x \sec ^{2} x d x$, let the variable $u$ in the Parts Formula be $x$, and let $d v$ be $\sec ^{2} x d x$. Then $d u=d x$ and $v=\tan x$, so
\[
\begin{aligned}
\displaystyle\int x \sec ^{2} x d x & =x \tan x-\displaystyle\int \tan x d x 
& =x \tan x+\ln |\cos x|+C
\end{aligned}
\]
\end{solution}



\newpage





\question
\[
\displaystyle\int \dfrac{2 x+1}{4+x^{2}} \,dx=
\]
\begin{choices}
  \choice $\ln \left(x^{2}+4\right)+C$
  \choice $\ln \left(x^{2}+4\right)+\tan ^{-1} \dfrac{x}{2}+C$
  \choice $\dfrac{1}{2} \tan ^{-1} \dfrac{x}{2}+C$
  \CorrectChoice $\ln \left(x^{2}+4\right)+\dfrac{1}{2} \tan ^{-1} \dfrac{x}{2}+C$
  \choice none of these
\end{choices}
\begin{solution}
The integral is equivalent to $\displaystyle\int \dfrac{2 x}{4+x^{2}} d x+\displaystyle\int \dfrac{1}{4+x^{2}} d x$. Use formula (4) on the first integral and (18) on the second.
\end{solution}

\question
\[
\displaystyle\int \dfrac{1-x}{\sqrt{1-x^{2}}} \,dx=
\]
\begin{choices}
  \choice $\sqrt{1-x^{2}}+C$
  \choice $\sin ^{-1} x+C$
  \choice $\dfrac{1}{2} \ln \sqrt{1-x^{2}}+C$
  \CorrectChoice $\sin ^{-1} x+\sqrt{1-x^{2}}+C$
  \choice $\sin ^{-1} x+\dfrac{1}{2} \ln \sqrt{1-x^{2}}+C$
\end{choices}
\begin{solution}
The integral is equivalent to $\displaystyle\int \dfrac{1}{\sqrt{1-x^{2}}} d x-\displaystyle\int \dfrac{x}{\sqrt{1-x^{2}}} d x$. Use formula (17) on the first integral. Rewrite the second integral as $-\dfrac{1}{2} \displaystyle\int\left(1-x^{2}\right)^{-\dfrac{1}{2}}(-2 x) d x$, and use (3).
\end{solution}

\question
\[
\displaystyle\int \dfrac{2 x-1}{\sqrt{4 x-4 x^{2}}} \,dx=
\]
\begin{choices}
  \choice $4 \ln \sqrt{4 x-4 x^{2}}+C$
  \choice $\sin ^{-1}(1-2 x)+C$
  \choice $\dfrac{1}{2} \sqrt{4 x-4 x^{2}}+C$
  \choice $-\dfrac{1}{4} \ln \left(4 x-4 x^{2}\right)+C$
  \CorrectChoice $-\dfrac{1}{2} \sqrt{4 x-4 x^{2}}+C$
\end{choices}
\begin{solution}
Rewrite: $-\dfrac{1}{4} \displaystyle\int\left(4 x-4 x^{2}\right)^{-\dfrac{1}{2}}(4-8 x) d x$.
\end{solution}


\newpage




\question
\[
\displaystyle\int \dfrac{e^{2 x}}{1+e^{x}} \,dx=
\]
\begin{choices}
  \choice $\tan ^{-1} e^{x}+C$
  \CorrectChoice $e^{x}-\ln \left(1+e^{x}\right)+C$
  \choice $e^{x}-x+\ln \left|1+e^{x}\right|+C$
  \choice $e^{x}+\dfrac{1}{\left(e^{x}+1\right)^{2}}+C$
  \choice none of these
\end{choices}
\begin{solution}
Hint: Divide, getting $\displaystyle\int\left[e^{x}-\dfrac{e^{x}}{1+e^{x}}\right] d x$.
\end{solution}

\question
\[
\displaystyle\int \dfrac{\cos \theta}{1+\sin ^{2} \theta} d \theta=
\]
\begin{choices}
  \choice $\sec \theta \tan \theta+C$
  \choice $\sin \theta-\csc \theta+C$
  \choice $\ln \left(1+\sin ^{2} \theta\right)+C$
  \CorrectChoice $\tan ^{-1}(\sin \theta)+C$
  \choice $-\dfrac{1}{\left(1+\sin ^{2} \theta\right)^{2}}+C$
\end{choices}
\begin{solution}
Letting $u=\sin \theta$ yields the integral $\displaystyle\int \dfrac{d u}{1+u^{2}}$. Use formula (18).
\end{solution}

\question
\[
\displaystyle\int \arctan x \,dx=
\]
\begin{choices}
  \choice $\operatorname{arc} \tan x+C$
  \choice $x \operatorname{arc} \tan x-\ln \left(1+x^{2}\right)+C$
  \choice $x \operatorname{arc} \tan x+\ln \left(1+x^{2}\right)+C$
  \choice $x \arctan x+\dfrac{1}{2} \ln \left(1+x^{2}\right)+C$
  \CorrectChoice $x \arctan x-\dfrac{1}{2} \ln \left(1+x^{2}\right)+C$
\end{choices}
\begin{solution}
Use integration by parts, letting $u=\arctan x$ and $d v=d x$. Then
\[
d u=\dfrac{d x}{1+x^{2}} \text { and } v=x
\]
The Parts Formula yields
\[
x \arctan x-\displaystyle\int \dfrac{x d x}{1+x^{2}} \quad \text { or } \quad x \arctan x-\dfrac{1}{2} \ln \left(1+x^{2}\right)+C
\]
\end{solution}


\newpage




\question
\[
\displaystyle\int \dfrac{d x}{1-e^{x}}=
\]
\begin{choices}
  \choice $-\ln \left|1-e^{x}\right|+C$
  \CorrectChoice $x-\ln \left|1-e^{x}\right|+C$
  \choice $\dfrac{1}{\left(1-e^{x}\right)^{2}}+C$
  \choice $e^{-x} \ln \left|1+e^{x}\right|+C$
  \choice none of these
\end{choices}
\begin{solution}
Hint: Note that


\[
\dfrac{1}{1-e^{x}}=\dfrac{1-e^{x}+e^{x}}{1-e^{x}}=1+\dfrac{e^{x}}{1-e^{x}}
\]

Or multiply the integrand by $\dfrac{e^{-x}}{e^{-x}}$, recognizing that the correct answer is equivalent to $-\ln \left|e^{-x}-1\right|$.
\end{solution}

\question
\[
\displaystyle\int \dfrac{(2-y)^{2}}{4 \sqrt{y}} \,dy=
\]
\begin{choices}
  \choice $\dfrac{1}{6}(2-y)^{3} \sqrt{y}+C$
  \choice $2 \sqrt{y}-\dfrac{2}{3} y^{3 / 2}+\dfrac{8}{5} y^{5 / 2}+C$
  \choice $\ln |y|-y+2 y^{2}+C$
  \CorrectChoice $2 y^{1 / 2}-\dfrac{2}{3} y^{3 / 2}+\dfrac{1}{10} y^{5 / 2}+C$
  \choice none of these
\end{choices}
\begin{solution}
(D) Hint: Expand the numerator and divide. Then integrate term by term.
\end{solution}

\question
\[
\displaystyle\int e^{2 \ln u} \,du=
\]
\begin{choices}
  \choice $\dfrac{1}{3} e^{u^{3}}+C$
  \choice $e^{u^{3} / 3}+C$
  \CorrectChoice $\dfrac{1}{3} u^{3}+C$
  \choice $\dfrac{2}{u} e^{2 \ln u}+C$
  \choice $e^{1+2 \ln u}+C$
\end{choices}
\begin{solution}
(C) Hint: Observe that $e^{2 \ln u}=u^{2}$.
\end{solution}


\newpage




\question
\[
\displaystyle\int \dfrac{d y}{y\left(1+\ln y^{2}\right)}=
\]
\begin{choices}
  \CorrectChoice $\dfrac{1}{2} \ln \left|1+\ln y^{2}\right|+C$
  \choice $-\dfrac{1}{\left(1+\ln y^{2}\right)^{2}}+C$
  \choice $\ln |y|+\dfrac{1}{2} \ln |\ln y|+C$
  \choice $\tan ^{-1}(\ln |y|)+C$
  \choice none of these
\end{choices}
\begin{solution}
(A) Let $u=1+\ln y^{2}=1+2 \ln |y|$; integrate $\dfrac{1}{2} \displaystyle\int \dfrac{\dfrac{2}{y} d y}{1+2 \ln |y|}$.
\end{solution}

\question
\[
\displaystyle\int(\tan \theta-1)^{2} d \theta=
\]
\begin{choices}
  \choice $\sec \theta+\theta+2 \ln |\cos \theta|+C$
  \CorrectChoice $\tan \theta+2 \ln |\cos \theta|+C$
  \choice $\tan \theta-2 \sec ^{2} \theta+C$
  \choice $\sec \theta+\theta-\tan ^{2} \theta+C$
  \choice $\tan \theta-2 \ln |\cos \theta|+C$
\end{choices}
\begin{solution}
(B) Hint: Expand and note that
\[
\displaystyle\int\left(\tan ^{2} \theta-2 \tan \theta+1\right) d \theta=\displaystyle\int \sec ^{2} \theta d \theta-2 \displaystyle\int \tan \theta d \theta
\]
Use formulas (9) and (7).
\end{solution}


\newpage




\question
\[
\displaystyle\int \dfrac{d \theta}{1+\sin \theta}=
\]
\begin{choices}
  \choice $\sec \theta-\tan \theta+C$
  \choice $\ln (1+\sin \theta)+C$
  \choice $\ln |\sec \theta+\tan \theta|+C$
  \choice $\theta+\ln |\csc \theta-\cot \theta|+C$
  \CorrectChoice none of these
\end{choices}
\begin{solution}
(E) Multiply by $\dfrac{1-\sin \theta}{1-\sin \theta}$. The correct answer is $\tan \theta-\sec \theta+C$.
\end{solution}

\question
A particle starting at rest at $t=0$ moves along a line so that its acceleration at time $t$ is $12 t \mathrm{ft} / \mathrm{sec}^{2}$. How much distance does the particle cover during the first 3 sec ?
\begin{choices}
  \choice$16$ ft
  \choice$32$ ft
  \choice$48$ ft
  \CorrectChoice$54$ ft
  \choice $108 \mathrm{ft}$
\end{choices}
\begin{solution}
(D) Note the initial conditions: when $t=0, v=0$ and $s=0$. Integrate twice: $v=6 t^{2}$ and $s=2 t^{3}$. Let $t=3$.
\end{solution}

\question
The equation of the curve whose slope at point $(x, y)$ is $x^{2}-2$ and which contains the point $(1,-3)$ is
\begin{choices}
  \choice $y=\dfrac{1}{3} x^{3}-2 x$
  \choice $y=2 x-1$
  \choice $y=\dfrac{1}{3} x^{3}-\dfrac{10}{3}$
  \CorrectChoice $y=\dfrac{1}{3} x^{3}-2 x-\dfrac{4}{3}$
  \choice $3 y=x^{3}-10$
\end{choices}
\begin{solution}
(D) Since $y^{\prime}=x^{2}-2, y=\dfrac{1}{3} x^{3}-2 x+C$. Replacing $x$ by 1 and $y$ by -3 yields $C=-\dfrac{4}{3}$.
\end{solution}


\newpage




\question
A particle moves along a line with acceleration $2+6 t$ at time $t$. When $t=0$, its velocity equals 3 and it is at position $s=2$. When $t=1$, it is at position $s=$
\begin{choices}
  \choice$2$
  \choice$5$
  \choice$6$
  \CorrectChoice$7$
  \choice$8$
\end{choices}
\begin{solution}
(D) When $t=0, v=3$ and $s=2$, so

\[
v=2 t+3 t^{2}+3 \quad \text { and } \quad s=t^{2}+t^{3}+3 t+2 .
\]

Let $t=1$.
\end{solution}

\question
Find the acceleration (in ft/ $\mathrm{sec}^{2}$ ) needed to bring a particle moving with a velocity of $75 \mathrm{ft} / \mathrm{sec}$ to a stop in 5 sec .
\begin{choices}
  \choice$-3$
  \choice$-6$
  \CorrectChoice$-15$
  \choice$-25$
  \choice$-30$
\end{choices}
\begin{solution}
(C) Let $\dfrac{d v}{d t}=a$; then
\begin{equation*}
v=a t+C . \tag{*}
\end{equation*}
Since $v=75$ when $t=0$, therefore $C=75$. Then (*) becomes
\[
v=a t+75
\]
so
\[
0=a t+75 \quad \text { and } \quad a=-15
\]
\end{solution}


\newpage




\question\label{calc05:bc-last}
\[
\displaystyle\int \dfrac{x^{2}}{x^{2}-1} \,dx=
\]
\begin{choices}
  \CorrectChoice $x+\dfrac{1}{2} \ln \left|\dfrac{x-1}{x+1}\right|+C$
  \choice $\ln \left|x^{2}-1\right|+C$
  \choice $x+\tan ^{-1} x+C$
  \choice $x+\dfrac{1}{2} \ln \left|\dfrac{x+1}{x-1}\right|+C$
  \choice $1+\dfrac{1}{2} \ln \left|\dfrac{x+1}{x-1}\right|+C$
\end{choices}
\begin{solution}
Divide to obtain $\displaystyle\int\left(1+\dfrac{1}{x^{2}-1}\right) d x$. Use partial fractions to get


\[
\dfrac{1}{x^{2}-1}=\dfrac{1}{2(x-1)}-\dfrac{1}{2(x+1)} .
\]
\end{solution}